\documentclass[12pt]{article}

\usepackage[spanish]{babel}
\usepackage[utf8]{inputenc}
\usepackage[T1]{fontenc}
\usepackage{amsmath,amssymb,amsthm}
\usepackage{geometry}
\geometry{margin=2.5cm}

\title{Estructuras Algebraicas}
\date{}

\begin{document}

\maketitle

\section{Estructuras algebraicas}

Las estructuras algebraicas que estudiaremos principalmente son los grupos, anillos y módulos.
Todas estas estructuras se engloban en las llamadas \emph{álgebras universales}.

\section{Operaciones \(n\)-arias}

Sea
\[
\mathbb{N}=\{0,1,2,3,\dots\}
\]
el conjunto de los números naturales, y sea \(A\) un conjunto.

Para cada \(n\in\mathbb{N}\), una \textbf{operación \(n\)-aria sobre \(A\)} es una función
\[
\omega:A^n\to A,
\]
donde
\[
A^n=\underbrace{A\times A\times\cdots\times A}_{n\text{ veces}}.
\]

Por convención:
\[
A^0 \text{ es un conjunto unitario}, \qquad A^1=A.
\]

\begin{center}
\begin{tabular}{c|c|c}
\(n\) & operación & tipo \\
\hline
0 & 0-aria & \(\{*\}\to A\) \\
1 & 1-aria & \(A\to A\) \\
2 & binaria & \(A\times A\to A\) \\
3 & ternaria & \(A\times A\times A\to A\)
\end{tabular}
\end{center}

\paragraph{Observación.}
Dar operaciones \(0\)-arias sobre \(A\) equivale a dar elementos de \(A\).

\section{Operaciones binarias}

Sea
\[
\omega:A\times A\to A.
\]

Diremos que \(\omega\) es \textbf{asociativa} si
\[
\omega(\omega(a,b),c)=\omega(a,\omega(b,c))
\qquad \forall a,b,c\in A.
\]

Usando la notación usual
\[
a\cdot b=\omega(a,b),
\]
la asociatividad se escribe como
\[
(a\cdot b)\cdot c=a\cdot (b\cdot c)
\qquad \forall a,b,c\in A.
\]

\section{Primeras estructuras}

\subsection{Magma}

Un conjunto \(A\) con una operación binaria se llama \textbf{magma}.

Se denota por
\[
(A,\cdot).
\]

\subsection{Semigrupo}

Un conjunto no vacío \(A\) junto con una operación binaria asociativa se llama \textbf{semigrupo}.

\subsection{Elemento neutro}

Sea \((A,\cdot)\) un magma.

Un elemento \(e\in A\) se llama \textbf{elemento neutro} si
\[
e\cdot a=a=a\cdot e
\qquad \forall a\in A.
\]

\subsection{Monoide}

Un \textbf{monoide} es un semigrupo con elemento neutro.

Se denota por
\[
(A,\cdot,e).
\]

\subsection{Magma conmutativo}

Un magma \(A\) es \textbf{conmutativo} o \textbf{abeliano} si
\[
a\cdot b=b\cdot a
\qquad \forall a,b\in A.
\]

\section{Propiedades básicas}

\subsection{Unicidad del elemento neutro}

\textbf{Proposición.}
En un monoide, el elemento neutro es único.

\textbf{Demostración.}
Sean \(e\) y \(e'\) dos elementos neutros. Entonces, como \(e\) es neutro,
\[
e\cdot e'=e',
\]
y como \(e'\) es neutro,
\[
e\cdot e'=e.
\]
Por lo tanto,
\[
e=e'.
\]
\qed

\subsection{Elemento inversible}

Sea \((A,\cdot,e)\) un magma con elemento neutro.

Un elemento \(a\in A\) se dice \textbf{invertible} si existe \(b\in A\) tal que
\[
ab=e=ba.
\]

En ese caso, \(b\) se llama el \textbf{inverso} de \(a\).

\subsection{Unicidad del inverso}

\textbf{Proposición.}
En un monoide, si un elemento tiene inverso, entonces dicho inverso es único.

\textbf{Demostración.}
Supongamos que \(b\) y \(b'\) son inversos de \(a\). Entonces
\[
ab=e=ba
\qquad\text{y}\qquad
ab'=e=b'a.
\]
Así,
\[
b=b\cdot e=b(ab')=(ba)b'=e\,b'=b'.
\]
Por lo tanto, el inverso es único.
\qed

\paragraph{Notación.}
Si \(b\) es el inverso de \(a\), se escribe
\[
a^{-1}=b.
\]

\section{Grupo}

Un \textbf{grupo} es un monoide en el que todo elemento es invertible.

Se denota por
\[
(G,\cdot,{}^{-1},e).
\]

Aquí:
\begin{itemize}
\item \(\cdot\) es una operación binaria;
\item \({}^{-1}\) es una operación 1-aria;
\item \(e\) puede verse como una operación 0-aria.
\end{itemize}

\section{Ejemplos ordenados}

\subsection{Ejemplo de monoide no conmutativo}

Consideremos
\[
M_{2\times 2}(\mathbb{Z}),
\]
el conjunto de matrices cuadradas \(2\times 2\) con coeficientes enteros.

Con el producto usual de matrices, la identidad
\[
I=\begin{pmatrix}1&0\\0&1\end{pmatrix}
\]
es el elemento neutro. Por tanto,
\[
\bigl(M_{2\times 2}(\mathbb{Z}),\cdot,I\bigr)
\]
es un monoide.

\subsection{Ejemplos de monoides abelianos}

\[
(\mathbb{N},+,0)
\qquad\text{y}\qquad
(\mathbb{N},\cdot,1)
\]
son monoides abelianos distintos.

\subsection{Ejemplo de operación binaria no asociativa}

Sea
\[
A=M_{2\times 2}(\mathbb{Z})
\]
y definamos
\[
\omega:A\times A\to A,
\qquad
\omega(M,N)=MN-NM.
\]

Esta operación \textbf{no es asociativa}.

\textbf{Demostración.}
Tomemos
\[
A_1=\begin{pmatrix}1&0\\0&0\end{pmatrix},
\qquad
A_2=\begin{pmatrix}0&1\\0&0\end{pmatrix},
\qquad
A_3=\begin{pmatrix}0&0\\1&0\end{pmatrix}.
\]

Primero calculamos
\[
\omega(A_1,A_2)=A_1A_2-A_2A_1
=
\begin{pmatrix}0&1\\0&0\end{pmatrix}
-
\begin{pmatrix}0&0\\0&0\end{pmatrix}
=
\begin{pmatrix}0&1\\0&0\end{pmatrix}.
\]
Entonces
\[
\omega(\omega(A_1,A_2),A_3)
=
\omega\!\left(
\begin{pmatrix}0&1\\0&0\end{pmatrix},
\begin{pmatrix}0&0\\1&0\end{pmatrix}
\right).
\]
Ahora,
\[
\begin{pmatrix}0&1\\0&0\end{pmatrix}
\begin{pmatrix}0&0\\1&0\end{pmatrix}
=
\begin{pmatrix}1&0\\0&0\end{pmatrix},
\qquad
\begin{pmatrix}0&0\\1&0\end{pmatrix}
\begin{pmatrix}0&1\\0&0\end{pmatrix}
=
\begin{pmatrix}0&0\\0&1\end{pmatrix}.
\]
Por tanto,
\[
\omega(\omega(A_1,A_2),A_3)
=
\begin{pmatrix}1&0\\0&-1\end{pmatrix}.
\]

Por otro lado,
\[
\omega(A_2,A_3)=A_2A_3-A_3A_2
=
\begin{pmatrix}1&0\\0&0\end{pmatrix}
-
\begin{pmatrix}0&0\\0&1\end{pmatrix}
=
\begin{pmatrix}1&0\\0&-1\end{pmatrix}.
\]
Entonces
\[
\omega(A_1,\omega(A_2,A_3))
=
\omega\!\left(
\begin{pmatrix}1&0\\0&0\end{pmatrix},
\begin{pmatrix}1&0\\0&-1\end{pmatrix}
\right).
\]
Pero estas dos matrices conmutan, así que
\[
\omega(A_1,\omega(A_2,A_3))=0.
\]

En consecuencia,
\[
\omega(\omega(A_1,A_2),A_3)\neq \omega(A_1,\omega(A_2,A_3)).
\]
Luego \(\omega\) no es asociativa.
\qed

\subsection{Ejemplo de no invertibilidad en \(M_{2\times 2}(\mathbb{Z})\)}

En \(\bigl(M_{2\times 2}(\mathbb{Z}),\cdot,I\bigr)\), la matriz \(2I\) no es invertible.

\textbf{Demostración.}
Se tiene
\[
\det(2I)=4\neq 0.
\]
Sin embargo, si \(2I\) fuera invertible en \(M_{2\times 2}(\mathbb{Z})\), existiría
\[
N\in M_{2\times 2}(\mathbb{Z})
\]
tal que
\[
(2I)N=I.
\]
De aquí se deduce
\[
2N=I,
\qquad\text{es decir}\qquad
N=\frac12 I
=
\begin{pmatrix}
\frac12 & 0\\
0 & \frac12
\end{pmatrix}.
\]
Pero esta matriz no pertenece a \(M_{2\times 2}(\mathbb{Z})\), contradicción.

Por lo tanto, \(2I\) no es invertible en \(M_{2\times 2}(\mathbb{Z})\).
\qed

\paragraph{Observación.}
Todo elemento de \(\bigl(M_{2\times 2}(\mathbb{Z}),+\bigr)\) es invertible, pues el inverso aditivo de \(M\) es \(-M\).

Por tanto,
\[
\bigl(M_{2\times 2}(\mathbb{Z}),+\bigr)
\]
es un grupo, mientras que
\[
\bigl(M_{2\times 2}(\mathbb{Z}),\cdot\bigr)
\]
no es un grupo.

\subsection{Grupo de biyecciones}

Sea \(X\) un conjunto no vacío y definamos
\[
\operatorname{Bij}(X)=\{f:X\to X:\ f \text{ es biyectiva}\}.
\]
Entonces,
\[
\bigl(\operatorname{Bij}(X),\circ\bigr)
\]
es un grupo.

\subsection{Grupo simétrico}

Si \(|X|=n\), entonces
\[
\Sigma_n=\bigl(\operatorname{Bij}(X),\circ\bigr)
\]
se llama el \textbf{grupo de permutaciones de \(n\) letras}.

En particular,
\[
|\Sigma_n|=n!.
\]

\section{Álgebras universales}

Sea \(\Omega\) un conjunto y sea
\[
a:\Omega\to\mathbb{N}
\]
una función.

Para cada \(n\in\mathbb{N}\), definimos
\[
\Omega(n)=\{\omega\in\Omega:\ a(\omega)=n\}.
\]

\textbf{Definición.}
Una \(\Omega\)-álgebra, o álgebra universal, es una terna
\[
(A,\Omega,\varphi),
\]
donde \(A\) es un conjunto y \(\varphi\) es una familia de funciones
\[
\varphi=\{\varphi_n\}_{n\in\mathbb{N}},
\qquad
\varphi_n:\Omega(n)\to \operatorname{Func}(A^n,A).
\]

Equivalentemente, para cada \(\omega\in\Omega(n)\),
\[
\varphi_n(\omega):A^n\to A
\]
es una operación \(n\)-aria sobre \(A\).

\paragraph{Abuso de notación.}
Muchas veces
\[
(A,\Omega,\varphi)
\]
se escribe simplemente como
\[
(A,\varphi)
\]
o incluso indicando solo las operaciones relevantes.

\section{Interpretación de estructuras como álgebras universales}

\subsection{Magma}

Un magma puede verse como
\[
(A,\omega),
\]
donde
\[
\omega:A\times A\to A
\]
es una operación binaria.

Usualmente se escribe
\[
(A,\cdot).
\]

\subsection{Monoide}

Un monoide puede verse como
\[
(A,\cdot,e),
\]
donde:
\begin{itemize}
\item \(\cdot\) es una operación binaria;
\item \(e\) es una operación 0-aria, es decir, una constante del conjunto.
\end{itemize}

\subsection{Grupo}

Un grupo puede verse como
\[
(G,\cdot,{}^{-1},e),
\]
donde:
\begin{itemize}
\item \(\cdot\) es binaria;
\item \({}^{-1}\) es 1-aria;
\item \(e\) es 0-aria.
\end{itemize}

\section{Ejemplo: espacios vectoriales}

Sea \(K\) un cuerpo. Los \(K\)-espacios vectoriales pueden interpretarse como álgebras universales.

Sea \(V\) un \(K\)-espacio vectorial. Entonces tenemos:

\[
+ : V\times V\to V,
\qquad
(u,v)\mapsto u+v,
\]
que es una operación binaria;

para cada \(\alpha\in K\), la aplicación
\[
w_\alpha:V\to V,
\qquad
v\mapsto \alpha v,
\]
que es una operación 1-aria;

y además el elemento
\[
0\in V,
\]
que puede verse como una operación 0-aria.

Por tanto, un espacio vectorial puede escribirse como una estructura del tipo
\[
\bigl(V,+,\{w_\alpha\}_{\alpha\in K},0\bigr),
\]
donde:
\begin{itemize}
\item \(+\) es binaria;
\item cada \(w_\alpha\) es 1-aria;
\item \(0\) es 0-aria y corresponde al neutro aditivo.
\end{itemize}

\end{document}